% ==============================================================================
% 03_methodology.tex – Chapter 3
% ==============================================================================
\chapter{Methodology}
\label{ch:methodik}

This chapter gives an overview of the procedure: data
(\cref{sec:meth_daten,sec:meth_aufnahme}), pre-processing
(\cref{sec:meth_vorverarbeitung}), model and training (\cref{sec:meth_modell}),
evaluation (\cref{sec:meth_evaluation}) and implementation
(\cref{sec:meth_implementierung}). The exact settings of the experiments are
given in \cref{sec:setup}.

\section{Datasets}
\label{sec:meth_daten}
The thesis uses two data sources:
\begin{itemize}
  \item \textbf{GazeCapture}~\cite{krafka2016itracker} for the base models. The
        full dataset has \num{2445504} frames of \num{1474} persons. This
        thesis uses a publicly available subset from
        Kaggle~\cite{kaggle_gazecapture}. For every person, the frames and the
        corresponding target points on the screen are read from the metadata.
        Persons with incomplete metadata are skipped. This gives \num{347085}
        labelled frames of 262 persons on iPhones and iPads. There are two splits, each 80/20:
        \code{norm\_subject} (\SI{80}{\percent} of the persons in training,
        \SI{20}{\percent} in test; no person in both) and \code{norm\_random} (random by
        groups of person and target point; the same person can appear in
        training and test). For the experiments, \num{1000} training and
        \num{200} test frames are drawn at random, plus 200 validation frames
        of a further person.
  \item \textbf{Own recordings} for personalization. Person~A (the author)
        recorded 27 calibration sessions with two laptop webcams; 18 of them are
        evaluated. For comparison, a second person (person~B) recorded two
        sessions (\cref{sec:meth_aufnahme}).
\end{itemize}

\section{Recording the Personalization Data}
\label{sec:meth_aufnahme}
A recording program shows a red point on the screen, and the person follows it
with the eyes. At the same time, the webcam records images. For each image, the
position of the point is stored, with 6 labels per second. The positions are directly
normalized to $[0,1]^2$. Three generations of the program are used:
\begin{itemize}
  \item \textbf{Gen~1} (sessions 01--04): the point moves smoothly and without
        jumps across the whole screen.
  \item \textbf{Gen~2} (sessions 05--09): the point jumps every 1.4 to 2.7
        seconds to a random new place. The recordings are short (about
        \SI{30}{\second}).
  \item \textbf{Gen~4} (sessions 19--27): the recording consists of segments of
        \SI{5}{\second}. Each segment starts with \SI{0.5}{\second}
        \code{WAIT} at a new start point; then the point moves for
        \SI{4.5}{\second} (\code{MOVE}). \SI{90}{\percent} of the start points
        lie in the corners. The program stores the phase for every frame.
\end{itemize}
Gen~3 (sessions 10--18) was discarded because the model could not learn
anything from it. The sessions also differ in light, background, glasses, length
(90 or \SI{120}{\second} in Gen~4) and device. Two laptops, each with a built-in
webcam, were used: \emph{device~A}, a Microsoft Surface Pro~9 (screen area
\num{1430}~$\times$~\num{880} pixels), for sessions 01--03 and 25--27, and
\emph{device~B}, an HP laptop (\num{1870}~$\times$~\num{1000} pixels), for
sessions 05--09 and 19--24. For session 04 the device was not logged. On both devices the
frames are stored at \num{640}~$\times$~\num{480} pixels.
\Cref{tab:session_conditions} shows all conditions.

For comparison, a second person (person~B) recorded two sessions with Gen~4
(\SI{120}{\second}, \SI{90}{\percent} corners, without glasses, bright):
session~28 with device~A in front of a calm white background and session~29 with
device~B in front of a busy background with a window. Session~28 uses a slightly
smaller screen area (\num{1390}~$\times$~\num{880} pixels); the normalization
to $[0,1]^2$ compensates for this. The sessions of person~B are evaluated
separately (\cref{sec:person_b}) and are not part of the means over the 18
sessions of person~A.

\section{Pre-Processing}
\label{sec:meth_vorverarbeitung}
\begin{itemize}
  \item \textbf{Targets:} the gaze targets of GazeCapture are divided by the
        width and height of the respective screen. So all targets lie in the
        unit square $[0,1]^2$, independent of the device
        (\cref{sec:normalization}).
  \item \textbf{Images:} the whole camera image is resized to
        $224\times224$ pixels and normalized with the ImageNet means and
        standard deviations. There is no face or eye crop.
  \item \textbf{Cleaned labels:} for Gen~4, the \code{WAIT} frames can be
        removed, because the phase of every frame is logged
        (\cref{sec:cleaning}).
\end{itemize}

\section{Model and Training}
\label{sec:meth_modell}
The model has two parts:
\begin{enumerate}
  \item a \textbf{ResNet-18}~\cite{he2016resnet}, pretrained on ImageNet
        as feature extractor. It gives
        512 features per image.
  \item a \textbf{regression head}: an MLP of four fully connected layers
        $512 \to 512 \to 256 \to 128 \to 2$ with ReLU, dropout $p=0.2$ after
        the second layer and an output activation (default: sigmoid). It
        computes the position $(x,y)$ from the features.
\end{enumerate}
Training uses L1 loss, AdamW (learning rate $10^{-4}$,
weight decay $10^{-5}$), batch size 32 and early stopping (stop after three
epochs without improvement, at most 500 epochs). The procedure has two stages:
\begin{enumerate}
  \item \textbf{Base model (E0):} the head is trained on GazeCapture. Of five seeds, the run with the lowest
        validation error is kept.
  \item \textbf{Personalization (E1--E3):} the head of the base model is trained
        further with the frames of one own session.
        The backbone stays frozen, including the BatchNorm statistics.
\end{enumerate}
Because the backbone is frozen, its features are computed only once and stored. After that, only the head is trained. A run therefore takes
only seconds.

\section{Evaluation}
\label{sec:meth_evaluation}
\begin{itemize}
  \item \textbf{Error measure:} diagonal error $\Ediag$, that is, the Euclidean
        RMSE in percent of the screen diagonal (\cref{eq:ediag}).
  \item \textbf{Baselines:} the constant baseline always predicts the mean of
        the training targets; the random baseline guesses at random. Both are
        computed on the same test data as the model (\cref{sec:baselines}).
  \item \textbf{Statistics:} for GazeCapture there are five repeats per setting
        with a \SI{98}{\percent} confidence interval (Student's $t$) and the
        Mann-Whitney $U$ test. For the sessions, I use the Wilcoxon signed-rank
        test (paired) and the Kruskal-Wallis test (comparison of groups).
  \item \textbf{Splitting the session data:} random 70/15/15, sequential
        (temporal blocks with a gap of one second) and two videos (training on
        one session, test on another) (\cref{sec:loading}).
  \item \textbf{Spatial analysis:} errors in the centre, at the edges and in the
        corners of the screen (\cref{sec:partition}).
  \item \textbf{Temporal smoothing (E4):} a moving average over the last $k$
        predictions is applied afterwards to the stored test predictions.
        Accuracy ($\Ediag$) and stability (RMS-S2S) are measured
        (\cref{sec:smoothing}).
\end{itemize}

\section{Implementation}
\label{sec:meth_implementierung}
The code is written in Python with PyTorch and torchvision. A short overview of
the code structure and the commands to repeat the experiments are given in the
appendix (\cref{sec:anhang_code,sec:anhang_befehle}).
All experiments run without a GPU on the CPU of the HP laptop
(device~B): AMD Ryzen~7 3700U, about \SI{13}{\giga\byte} usable memory,
Ubuntu~26.04.1~LTS, Python~3.12.3, PyTorch~2.12.1 and torchvision~0.27.1
(CPU versions). The code and the web demo are
available on GitHub: \url{https://github.com/RafeiKAr/gaze-estimation}. The base model of the demo is
available on Hugging Face: \url{https://huggingface.co/RafeiKAr/demo_model}. A copy
is stored in the BHT cloud:
\url{https://cloud.bht-berlin.de/index.php/s/P4GkQ9J8zWHbdf9}. The frames
and videos of the own recordings show the faces of the two persons. They are therefore not
public, but stored in a private dataset on Hugging Face.


